\documentclass[10pt,a4paper]{amsart}
\usepackage[margin=19mm]{geometry}
\usepackage{amsmath,amssymb,mathtools}
\usepackage[scr=rsfs,cal=euler]{mathalfa}
\usepackage{xfrac}
\usepackage[unicode,hidelinks,bookmarks=false]{hyperref}
\newcommand{\ie}{i.e.\ }
\newcommand{\cf}{cf.\ }
\newcommand{\resp}{respectively\ }
\newcommand{\Subsection}{\subsection}
\providecommand{\nobreakdash}{\nobreak\hbox{-}}
\setlength{\emergencystretch}{2em}
\multlinegap0pt
\usepackage{tikz}
\usetikzlibrary{cd}
\usepackage{amssymb}
\usepackage{xcolor}
\usepackage{tikz}
\newtheorem{proposition}{Proposition}
\newtheorem{remark}{Remark}
\newcommand{\PrCorr}{\mathsf{Corr}_{\mathrm{pr}}}
\newcommand{\PrGrCorr}{\mathsf{GrCorr}_{\mathrm{pr}}}
\DeclareMathOperator{\Res}{Res}
\newcommand{\Span}{\mathsf{Span}}
\newcommand{\TDLC}{\mathrm{TDLC}}
\newcommand{\id}{\mathrm{id}}
\newcommand{\pret}{\mathrm{pr,\acute{e}t}}
\title{Chern character for torsion-free ample groupoids}
\author{Valerio Proietti, Makoto Yamashita}
\date{Source: arXiv:2509.07663v2 (CC BY 4.0)}
\begin{document}
\maketitle

\noindent\emph{Source and licence.} Chern character for torsion-free ample groupoids. Version arXiv:2509.07663v2 (CC BY 4.0), distributed by arXiv under the Creative Commons Attribution 4.0 International licence, http://creativecommons.org/licenses/by/4.0/, as recorded in the arXiv licence metadata for this exact version and shown on the abstract page https://arxiv.org/abs/2509.07663v2. Full licence notice: \texttt{licenses/arxiv-2509.07663-CC-BY-4.0.txt}.\par\smallskip

\noindent\textit{Editorial scope.} Counted unit: the article's proposition prop:span-corr-comm (source0.tex lines 1087--1100). The article's complete original proof of it is delivered with it (source0.tex lines 1150--1165, one \texttt{proof} environment of 16 lines, which the article places away from the statement). No mathematics is written, repaired or re-derived: the statement and the proof are the article's own lines, in the article's own notation, and no retained line is substituted or re-flowed.  Retained uncounted as the source-local context the proof needs: lines 1117--1140.  Nothing was omitted from any retained span, so the package carries no omission record.  No retained line contains a citation, so the wrapper carries no bibliography and no bibliography entry.  The retained lines contain the article's own diagram code, which the wrapper typesets with the standard tikz packages; no image file is delivered and the package declares no asset dependency.  Editorial formatting only, in addition: an AMS \texttt{amsart} preamble that restates the article's own declarations and macros used by the retained lines, namely the environments \texttt{proposition}, \texttt{remark} on separate counters, together with the standard packages those lines need.  The retained environments print in this document as Proposition 1, Remark 1 in order of appearance, whereas the article prints the same items with the numbering of its own full document.  The complete native source of the article as distributed in the arXiv e-print for this exact version is bundled unchanged as \texttt{sources/184758/source0.tex}.
\begin{remark}\label{rem:diagj}
It is not hard to verify that the diagram
\begin{equation}\label{eq:diagj}
\begin{tikzcd}
\Span_{\pret}(\TDLC^G) \arrow[r, "E^G"] \arrow[d, "G\ltimes -"] & \PrCorr^G \arrow[d, "- \rtimes G"] \\
\PrGrCorr \arrow[r, "E"]                 & \PrCorr               
\end{tikzcd}
\end{equation}
commutes up to natural equivalence.
Indeed, if we denote by $\mathfrak E$ the upper semicontinuous bundle associated to a Hilbert $G$-module $E$, then $E\rtimes G$ is a completion of $\Gamma_c(G;r^*\mathfrak{E})$ with inner product
\begin{equation}\label{eq:innprod}
(\xi|\zeta) (\gamma)= \sum_{s(\eta)=r(\gamma)} \eta^{-1}\cdot ({\xi}(\eta)|\zeta(\eta\gamma)).
\end{equation}
For a $G$-span $X \leftarrow Z \to Y$, we have that $E^G(Z)$ is a completion of $\Gamma_c(Z)$, so that $\Gamma_c(G;r^*\mathfrak{E})$ is a completion of $\Gamma_c(Z\times_{\rho,r} G)$ as prescribed by $E(G\ltimes \Omega^G_Z)$. The inner product in~\eqref{eq:innprod} corresponds to
\begin{equation*}%\label{eq:innprof}
(\xi|\zeta) (\gamma,y)= \sum_{s(\eta)=r(\gamma)} \sum_{f(z)=\eta\cdot y}(\bar{\xi}(z,\eta)|\zeta(z,\eta\gamma)),
\end{equation*}
which can be rewritten as 
\begin{equation}\label{eq:innprod2}
(\xi|\zeta) (\gamma,y)= \sum_{d(x)=y} \bar{\xi}(x)|\zeta(x\cdot (\gamma,y)),
\end{equation}
where $d\colon Z\times_{\rho,r} G\to Y$ is the right anchor map defined above for $G\ltimes \Omega^G_Z$.
Equation~\eqref{eq:innprod2} is the inner product as prescribed by $E(G\ltimes \Omega^G_Z)$, thus the inner products match up. The verification of compatibility for left and right action is proved just as easily, hence passing to completion we have that diagram~\eqref{eq:diagj} does indeed commute up to isomorphism of C$^*$-correspondences.
\end{remark}

\begin{proposition}\label{prop:span-corr-comm}
The following diagram commutes up to $2$-categorical equivalence.
\[
\begin{tikzcd}
\Span_{\pret}(\TDLC^G) \arrow[r, "E^G"] \arrow[d, "G\backslash(-)"] & \PrCorr^G \arrow[d, "- \rtimes G"] \\
\Span_{\pret}(\TDLC) \arrow[r, "E"]                 & \PrCorr               
\end{tikzcd}
\]
More precisely, given  $X \leftarrow Z \to Y$ in $\Span_{\pret}(\TDLC^G)$, there is a natural isomorphism between 
\[
E^G_{q_Y} \circ  E^G(Z)\rtimes G \circ (E^G_{q_X})^{\mathrm{op}}\quad \mathrm{ and } \quad E(G\backslash Z)
\]
as correspondences $C_0(G\backslash X)\to C_0(G\backslash Y)$.
\end{proposition} 

\begin{proof}[Proof of Proposition~\ref{prop:span-corr-comm}]
The squares in the diagram
\[
\begin{tikzcd}
X  \arrow[d, "q_X"] & Z \arrow[l, "g"'] \arrow[r, "f"] \arrow[d, "q_Z"] & Y \arrow[d, "q_Y"] \\
G\backslash X  & G\backslash Z  \arrow[l, "\bar{g}"'] \arrow[r, "\bar{f}"]                 & G\backslash Y               
\end{tikzcd}
\]
commute.
Note that there is an isomorphism $\Omega^G_Z\cong \Omega^G_f\circ \Omega_G^g$, where $\Omega_G^g\colon X\to Z$ is given by $Z$ with left anchor map $g$ and right anchor map $\id_Z$. By functoriality $ G\ltimes \Omega^G_Z$ is isomorphic to $G\ltimes \Omega^G_f\circ G\ltimes \Omega_G^g$. 
We have an isomorphism $\Omega^{G}_{q_Y}\circ (G\ltimes \Omega^G_f) \cong \Omega^G_{q_Y f} \cong \Omega^G_{\bar{f} q_Z}\cong \Omega_{\bar{f}}\circ \Omega^G_{q_Z}$. Now we have an isomorphism $\Omega^G_{q_Z}\circ (G\ltimes \Omega^g_G)\cong \Res^G_{G\ltimes X}(\Omega^G_{q_Z})$, where the latter correspondence is just $\Omega^G_{q_Z}$ with left action induced along the morphism $G\ltimes X \to G$ with anchor map $g^{(0)}$. Consider the correspondence $\Omega^{\bar{g}}\circ \Omega^G_{q_X}$. Its underlying topological space is the fiber product, denoted $Z_0$, over $X \to G\backslash X \leftarrow G\backslash Z$. By the universal property of pullbacks, we get a map $Z\to Z_0$. This map is compatible with the left and right actions, giving us a map $\Res^G_{G\ltimes X}(\Omega^G_{q_Z})\to \Omega^{\bar{g}}\circ \Omega^G_{q_X}$. This map is an isomorphism because morphisms between principal $G$-bundles over the same base are necessarily isomorphisms. Overall, we have
 \[
\Omega^{G}_{q_Y}\circ (G\ltimes \Omega^G_Z) \cong \Omega^{G}_{q_Y}\circ (G\ltimes \Omega^G_f)\circ (G\ltimes \Omega^g_G) \cong \Omega_{\bar{f}}\circ \Omega^{\bar{g}}\circ \Omega^G_{q_X}\cong \Omega_{G\backslash Z}\circ \Omega^G_{q_X}.
\]
The claim follows applying $E^G$ and using that $E^G(-)\rtimes G \cong E(- \rtimes G)$ (see Remark~\ref{rem:diagj}).
\end{proof}

\end{document}
